\section{Experimental protocol}

\subsection{Data, task construction, and protocol boundary}

We evaluate ICEWS14s with the original AnoT train/validation/test split files and anomaly-generation protocol. The source data are represented as integer quadruples and grouped by discrete timestamps. For each evaluation timestamp, normal facts are the observed split facts. Conceptual negatives are generated by replacing a compatible entity or relation with one that is not supported by the training graph; temporal negatives retain $(s,r,o)$ but place an early training fact at a later timestamp; and missing positives are sampled held-out facts, while their negatives combine time and conceptual corruptions. The generator seeds Python and NumPy with $42$, uses conceptual and missing ratios of $0.6$, a time ratio of $0.3$, and the ICEWS14s temporal span of $200$. Its temporal candidate pool is restricted to the earliest ten training timestamps and excludes facts known in validation or test, matching the original early-candidate construction.

This protocol produces labels for evaluation; it does not make the test labels available to the proposed miner. The contrastive time module uses the same construction \emph{inside training} to define a reference distribution for rule discovery. It does not use validation or test anomaly labels during training. Validation contains the only labels used for candidate selection, feature selection, score weights, and thresholds. Test scoring then consumes frozen rule files with graph updates disabled. The code-level guard rejects a training path that reads validation/test artifacts except as a known-fact collision filter, rejects candidate selection outside the validation split, and rejects test evaluation with updates or an unfiltered candidate file.

\begin{table}[t]
\centering
\caption{Key implementation settings for the reported ICEWS14s configuration. All temporal-rule settings are specified in \texttt{contrastive\_conflict\_mining.py} and \texttt{contrastive\_conflict\_selection.py}.}
\label{tab:implementation}
\small
\begin{tabular}{@{}p{0.55\linewidth}p{0.27\linewidth}@{}}
\toprule
Setting & Value \\
\midrule
Random seed & 42 \\
Original AnoT temporal span / early timestamps & 200 / 10 \\
Training displaced-reference ratio & 0.30 \\
Gap boundaries & 0, 30, 90, 180, 365 \\
Training minimum anomaly / total hits & 10 / 20 \\
Validation minimum lift / AUC & 1.25 / 0.505 \\
Posterior direction / credibility & 0.90 / 0.95 \\
Maximum temporal candidates / selected rules & 2,000 / 500 \\
Conceptual calibration candidates & linear, logistic, or GB (80 trees; depth 1--3; lr 0.02/0.05/0.1) \\
Manual conflicts / missing contrastive support & disabled / disabled \\
\bottomrule
\end{tabular}
\end{table}

\subsection{Baselines, calibration, and measures}

We compare \textbf{AnoT}, the reproduced original detector with validation-tuned thresholds, against \textbf{CARGO-AnoT}, the full error-aware configuration. CARGO-AnoT uses contrastive temporal rules for time, \texttt{novelty\_role} calibration for conceptual errors, and train-only missing support for missing errors. Manual relation-conflict rules are disabled because they did not survive validation in the time workflow. The optional missing-contrastive-support module is also disabled in the reported configuration because its validation behavior was unstable.

For explicit fusion, validation searches the configured feature-weight grid and selects the threshold that maximizes \fbeta. The conceptual \texttt{fair\_select} path avoids choosing a feature set on test data: it evaluates candidate linear and shallow gradient-boosted calibrators on a temporal holdout inside validation, selects one specification, then refits that specification on validation before scoring test. We report Precision, \fbeta, and ROC-AUC. Precision and \fbeta depend on the selected threshold, whereas AUC measures ranking independently of a single operating point. This distinction is central to interpreting the missing-fact results.

\subsection{Reproducibility and audit artifacts}

Each publication-mode evaluation emits the frozen configuration, source snapshot, selected-rule-file hashes, prediction rows, and a result-registry record. The registry also stores the protocol-guard outcomes, the command line, and the final metrics. Candidate, evidence, and validation artifacts carry their own input/output hashes. These artifacts do not replace an independent reproduction, but they make it possible to establish exactly which rule file and which settings produced a reported score. We use the archived registry values in this draft; before submission, the full table must be regenerated from one frozen archive and accompanied by the corresponding prediction files.

We report Precision, \fbeta, and AUC. \fbeta is the metric implemented by the evaluation code, with $\beta=0.5$. The full \method configuration uses contrastive temporal rules, conceptual role evidence, and train-only missing-support evidence. It disables manual relation-conflict rules and the unstable missing-contrastive-support module.

